% ==============================================================================
% estilo.tex — Preámbulo LaTeX común de mis PDF
% Lo carga estilo.yaml con include-in-header. Editá acá el aspecto de todos
% los documentos.
% ==============================================================================

% 0) Fuente principal con respaldo: emojis a color (Noto Color Emoji) y
%    símbolos (DejaVu Sans) cuando Latin Modern Roman no tiene el carácter.
%    Solo funciona con lualatex; con otro motor no se aplica.
\ifLuaTeX
  \directlua{
    luaotfload.add_fallback("emojis",
      {"Noto Color Emoji:mode=harf;", "DejaVu Sans:"})
  }
  \setmainfont{Latin Modern Roman}[RawFeature={fallback=emojis}]
\fi

% 1) Listas con más aire entre ítems (Pandoc las deja a 0 pt).
\renewcommand{\tightlist}{%
  \setlength{\itemsep}{0.55em}\setlength{\parskip}{0pt}}

% 2) Tablas con más aire entre filas (el valor por defecto es 1). Se aplica
%    solo dentro de las tablas: si se pusiera global, también ensancharía las
%    líneas de autor y fecha del título.
\AddToHook{env/longtable/begin}{\renewcommand{\arraystretch}{1.5}}

% 3) Cuadros "Idea clave / Ojo / Para el examen": barra azul a la
%    izquierda y fondo celeste pastel, con aire arriba y abajo.
%    Solo usa LaTeX base y color (sin paquetes extra).
\definecolor{barcolor}{RGB}{74,111,165}
\definecolor{fondocolor}{RGB}{232,241,251}
\newsavebox{\cuadrobox}
\renewenvironment{quote}{%
  \par\addvspace{10pt}%
  \begin{lrbox}{\cuadrobox}%
  \begin{minipage}[t]{\dimexpr\linewidth-24pt\relax}%
  \setlength{\parskip}{0.5em}}%
  {\par\end{minipage}\end{lrbox}%
  \par\noindent\hbox{%
    {\color{barcolor}\vrule width 2.5pt
      height \dimexpr\ht\cuadrobox+6pt\relax
      depth \dimexpr\dp\cuadrobox+6pt\relax}%
    {\setlength{\fboxsep}{0pt}%
     \colorbox{fondocolor}{\vrule width 0pt
      height \dimexpr\ht\cuadrobox+6pt\relax
      depth \dimexpr\dp\cuadrobox+6pt\relax
      \hskip 10pt\usebox{\cuadrobox}\hskip 10pt}}}%
  \par\addvspace{10pt}}

% 4) Bloques de código: las líneas largas se cortan en vez de salirse de la
%    hoja, y la línea que continúa empieza con una flechita (↪).
%    breakanywhere permite cortar también dentro de una palabra larga.
\usepackage{fvextra}
\fvset{breaklines, breakanywhere, breaknonspaceingroup}
\RecustomVerbatimEnvironment{verbatim}{Verbatim}{}
\ifdefined\Highlighting
  \RecustomVerbatimEnvironment{Highlighting}{Verbatim}{commandchars=\\\{\}}
\fi

% 5) Palabras largas sin espacios (URLs, hashes, rutas) en tablas, listas o
%    párrafos: si no entran en la línea, se cortan entre dos letras en vez de
%    salirse de la hoja. Solo se aplica a palabras de 20 caracteres o más
%    seguidos; las palabras normales no se tocan. Solo lualatex.
%    (Este bloque no puede llevar comentarios de Lua ni los símbolos de
%    porcentaje, numeral y virgulilla, porque TeX los interpretaría primero.)
\ifLuaTeX
  \directlua{
    local GLYPH = node.id("glyph")
    local KERN = node.id("kern")
    local PENALTY = node.id("penalty")
    local MATH = node.id("math")
    local MINIMO = 20
    local function cortar(head)
      local enmat = false
      local n = head
      while n do
        if n.id == MATH then
          enmat = (n.subtype == 0)
          n = n.next
        elseif n.id == GLYPH and not enmat then
          local largo = 0
          local m = n
          while m and (m.id == GLYPH or (m.id == KERN and m.subtype == 0)) do
            if m.id == GLYPH then largo = largo + 1 end
            m = m.next
          end
          if largo >= MINIMO then
            local c = n
            while c and not (c == m) do
              if c.id == GLYPH then
                local sig = c.next
                if sig and sig.id == KERN then sig = sig.next end
                if sig and not (sig == m) and sig.id == GLYPH then
                  local p = node.new(PENALTY)
                  p.penalty = 9000
                  head, c = node.insert_after(head, c, p)
                end
              end
              c = c.next
            end
          end
          n = m
        else
          n = n.next
        end
      end
      return head
    end
    luatexbase.add_to_callback("pre_linebreak_filter", cortar, "cortar-palabras-largas")
  }
\fi

% 6) \needspace{altura}: salta de página si no hay ese espacio libre
%    (definición del paquete needspace, copiada para no cargar nada).
\makeatletter
\providecommand{\needspace}[1]{\begingroup\setlength{\dimen@}{#1}%
  \vskip\z@\@plus\dimen@\penalty-100\vskip\z@\@plus-\dimen@
  \vskip\dimen@\penalty9999\vskip-\dimen@\vskip\z@skip\endgroup}
\makeatother
